\documentclass[11pt,compress,t,notes=noshow, xcolor=table]{beamer}
\usepackage[]{graphicx}\usepackage[]{color}
% maxwidth is the original width if it is less than linewidth
% otherwise use linewidth (to make sure the graphics do not exceed the margin)
\makeatletter
\def\maxwidth{ %
  \ifdim\Gin@nat@width>\linewidth
    \linewidth
  \else
    \Gin@nat@width
  \fi
}
\makeatother

\newcommand{\citebutton}[2]{%
\beamergotobutton{\href{#2}{#1}}%
}

\newcommand{\blu}[1]{\textcolor{blue}{#1}}
\newcommand{\org}[1]{\textcolor{orange}{#1}}
\newcommand{\ques}{\textbf{\textcolor{red}{Question:  }}}
\newcommand{\questionssofar}{\begin{frame}\frametitle{Any questions?}\end{frame}}

\newcommand\warning{%
 \makebox[1.4em][c]{%
 \makebox[0pt][c]{\raisebox{.1em}{\scriptsize!}}%
 \makebox[0pt][c]{\color{red}\normalsize$\bigtriangleup$}}}%

\definecolor{fgcolor}{rgb}{0.345, 0.345, 0.345}
\newcommand{\hlnum}[1]{\textcolor[rgb]{0.686,0.059,0.569}{#1}}%
\newcommand{\hlstr}[1]{\textcolor[rgb]{0.192,0.494,0.8}{#1}}%
\newcommand{\hlcom}[1]{\textcolor[rgb]{0.678,0.584,0.686}{\textit{#1}}}%
\newcommand{\hlopt}[1]{\textcolor[rgb]{0,0,0}{#1}}%
\newcommand{\hlstd}[1]{\textcolor[rgb]{0.345,0.345,0.345}{#1}}%
\newcommand{\hlkwa}[1]{\textcolor[rgb]{0.161,0.373,0.58}{\textbf{#1}}}%
\newcommand{\hlkwb}[1]{\textcolor[rgb]{0.69,0.353,0.396}{#1}}%
\newcommand{\hlkwc}[1]{\textcolor[rgb]{0.333,0.667,0.333}{#1}}%
\newcommand{\hlkwd}[1]{\textcolor[rgb]{0.737,0.353,0.396}{\textbf{#1}}}%
\let\hlipl\hlkwb

\usepackage{framed}
\makeatletter
\newenvironment{kframe}{%
 \def\at@end@of@kframe{}%
 \ifinner\ifhmode%
  \def\at@end@of@kframe{\end{minipage}}%
  \begin{minipage}{\columnwidth}%
 \fi\fi%
 \def\FrameCommand##1{\hskip\@totalleftmargin \hskip-\fboxsep
 \colorbox{shadecolor}{##1}\hskip-\fboxsep
     % There is no \\@totalrightmargin, so:
     \hskip-\linewidth \hskip-\@totalleftmargin \hskip\columnwidth}%
 \MakeFramed {\advance\hsize-\width
   \@totalleftmargin\z@ \linewidth\hsize
   \@setminipage}}%
 {\par\unskip\endMakeFramed%
 \at@end@of@kframe}
\makeatother

\definecolor{shadecolor}{rgb}{.97, .97, .97}
\definecolor{messagecolor}{rgb}{0, 0, 0}
\definecolor{warningcolor}{rgb}{1, 0, 1}
\definecolor{errorcolor}{rgb}{1, 0, 0}
\newenvironment{knitrout}{}{} % an empty environment to be redefined in TeX

\usepackage{alltt}
\newcommand{\SweaveOpts}[1]{}  % do not interfere with LaTeX
\newcommand{\SweaveInput}[1]{} % because they are not real TeX commands
\newcommand{\Sexpr}[1]{}       % will only be parsed by R
\newcommand{\xmark}{\ding{55}}%


\usepackage[english]{babel}
\usepackage[utf8]{inputenc}

\usepackage{dsfont}
\usepackage{verbatim}
\usepackage{amsmath}
\usepackage{amsfonts}
\usepackage{amssymb}
\usepackage{bm}
\usepackage{csquotes}
\usepackage{multirow}
\usepackage{longtable}
\usepackage{booktabs}
\usepackage{enumerate}
\usepackage[absolute,overlay]{textpos}
\usepackage{psfrag}
\usepackage{algorithm}
\usepackage{algpseudocode}
\usepackage{eqnarray}
\usepackage{arydshln}
\usepackage{tabularx}
\usepackage{placeins}
\usepackage{tikz}
\usepackage{setspace}
\usepackage{colortbl}
\usepackage{mathtools}
\usepackage{wrapfig}
\usepackage{bm}
\usepackage{amsmath}
\usepackage{pifont}

\usetikzlibrary{shapes.multipart,shapes,arrows,automata,positioning,calc,chains,trees, shadows}
\tikzset{
  %Define standard arrow tip
  >=stealth',
  %Define style for boxes
  punkt/.style={
    rectangle,
    rounded corners,
    draw=black, very thick,
    text width=6.5em,
    minimum height=2em,
    text centered},
  % Define arrow style
  pil/.style={
    ->,
    thick,
    shorten <=2pt,
    shorten >=2pt,}
}

\tikzstyle{vec}=[draw, rectangle, fill = white, minimum width=5mm, minimum height=1cm, inner sep = 2pt]

\usepackage{subfig}

% Defines macros and environments
\usepackage{../../style/lmu-lecture}


\let\code=\texttt
\let\proglang=\textsf

\setkeys{Gin}{width=0.9\textwidth}

\setbeamertemplate{frametitle}{\expandafter\uppercase\expandafter\insertframetitle}

\usepackage{bbm}
% basic latex stuff
\newcommand{\pkg}[1]{{\fontseries{b}\selectfont #1}} %fontstyle for R packages
\newcommand{\lz}{\vspace{0.5cm}} %vertical space
\newcommand{\dlz}{\vspace{1cm}} %double vertical space
\newcommand{\oneliner}[1] % Oneliner for important statements
{\begin{block}{}\begin{center}\begin{Large}#1\end{Large}\end{center}\end{block}}


%new environments
\newenvironment{vbframe}  %frame with breaks and verbatim
{
 \begin{frame}[containsverbatim,allowframebreaks]
}
{
\end{frame}
}

\newenvironment{vframe}  %frame with verbatim without breaks (to avoid numbering one slided frames)
{
 \begin{frame}[containsverbatim]
}
{
\end{frame}
}

\newenvironment{blocki}[1]   % itemize block
{
 \begin{block}{#1}\begin{itemize}
}
{
\end{itemize}\end{block}
}

\newenvironment{fragileframe}[2]{  %fragile frame with framebreaks
\begin{frame}[allowframebreaks, fragile, environment = fragileframe]
\frametitle{#1}
#2}
{\end{frame}}


\newcommand{\myframe}[2]{  %short for frame with framebreaks
\begin{frame}[allowframebreaks]
\frametitle{#1}
#2
\end{frame}}

\newcommand{\remark}[1]{
  \textbf{Remark:} #1
}


\newenvironment{deleteframe}
{
\begingroup
\usebackgroundtemplate{\includegraphics[width=\paperwidth,height=\paperheight]{../style/color/red.png}}
 \begin{frame}
}
{
\end{frame}
\endgroup
}
\newenvironment{simplifyframe}
{
\begingroup
\usebackgroundtemplate{\includegraphics[width=\paperwidth,height=\paperheight]{../style/color/yellow.png}}
 \begin{frame}
}
{
\end{frame}
\endgroup
}\newenvironment{draftframe}
{
\begingroup
\usebackgroundtemplate{\includegraphics[width=\paperwidth,height=\paperheight]{../style/color/green.jpg}}
 \begin{frame}
}
{
\end{frame}
\endgroup
}
% https://tex.stackexchange.com/a/261480: textcolor that works in mathmode
\makeatletter
\renewcommand*{\@textcolor}[3]{%
  \protect\leavevmode
  \begingroup
    \color#1{#2}#3%
  \endgroup
}
\makeatother





\input{../../latex-math/basic-math.tex}
\input{../../latex-math/basic-ml.tex}
\usepackage{colortbl}

\newcommand{\titlefigure}{figure/selfsup.jpg}
\newcommand{\learninggoals}{
\item Understand the difference to other learning paradigms
\item Learn to recognize self-supervision when you see it}

\title{Transfer Learning}
% \author{}
\institute{\href{https://slds-lmu.github.io/lecture_dl4nlp/}{slds-lmu.github.io/lecture\_dl4nlp}}
\date{}

\begin{document}
\lecturechapter{Self-Supervision}
\lecture{Deep Learning for NLP}

% ------------------------------------------------------------------------------

\begin{frame}{Definition}

	\vfill 
	
	\textit{Unsupervised Learning:}

	\begin{itemize}
		\item No labels attached to the data
		\item Learn patterns / clusters from the features only
	\end{itemize}
	
	\textit{Supervised Learning:}

	\begin{itemize}
		\item (Gold) Labels attached to the data
		\item Learn from the association between features and labels
	\end{itemize}
	
	\textbf{Self-Supervised Learning:}

	\begin{itemize}
		\item No \textit{external} labels attached to the data\\
					$\to$ Samples with suitable labels can be generated from the known structure of the data itself 
		\item \textit{Technically} supervised learning, but \textit{no labeling effort} $+$ simultan-\\eous ability to generate massive amounts of labeled data points
	\end{itemize}
	
	\vfill 
	
\end{frame}

% ------------------------------------------------------------------------------

\begin{vbframe}{Self-supervised objectives}

\vfill

\textbf{Recap: Language modeling}
	
	\begin{itemize}
		\item \textit{Training objective:} Given a context, predict the next word
	\end{itemize}
	
\begin{block}{Illustration (context size $=$ 2)}
\begin{tabular}{|c|c|c|cccccc|}
\hline
\cellcolor{blue!15}The & \cellcolor{blue!65}quick & brown & fox & jumps & over & the & lazy & dog \\
\hline
\end{tabular}\\
$\Rightarrow$ \quad (the, quick)
\begin{tabular}{|c|c|c|c|ccccc|}
\hline
\cellcolor{blue!15}The & \cellcolor{blue!15}quick & \cellcolor{blue!65}brown & fox & jumps & over & the & lazy & dog \\
\hline
\end{tabular}\\
$\Rightarrow$ \quad ([the, quick], brown)
\begin{tabular}{|c|c|c|c|ccccc|}
\hline
The & \cellcolor{blue!15}quick & \cellcolor{blue!15}brown & \cellcolor{blue!65}fox & jumps & over & the & lazy & dog \\
\hline
\end{tabular}\\
$\Rightarrow$ \quad ([quick, brown], fox)
\begin{tabular}{|c|c|c|c|ccccc|}
\hline
The & quick & \cellcolor{blue!15}brown & \cellcolor{blue!15}fox & \cellcolor{blue!65}jumps & over & the & lazy & dog \\
\hline
\end{tabular}\\
$\Rightarrow$ \quad ([brown, fox], jumps)
\end{block}
	
\vfill

\end{vbframe}

% ------------------------------------------------------------------------------

\begin{vbframe}{Self-supervised objectives}

\vfill

\textbf{Recap: Skip-gram}
	
	\begin{itemize}
		\item \textit{Training objective:} Given a word, predict the neighbouring words
    \item \textit{Generation of samples:} Sliding fixed-size window over the text
	\end{itemize}
	
\begin{block}{Illustration}
\begin{tabular}{|c|c|c|cccccc|}
\hline
\cellcolor{blue!15}The & \cellcolor{blue!65}quick & \cellcolor{blue!65}brown & fox & jumps & over & the & lazy & dog \\
\hline
\end{tabular}\\
$\Rightarrow$ \quad (the, quick); (the, brown)
\begin{tabular}{|c|c|c|c|ccccc|}
\hline
\cellcolor{blue!65}The & \cellcolor{blue!15}quick & \cellcolor{blue!65}brown & \cellcolor{blue!65}fox & jumps & over & the & lazy & dog \\
\hline
\end{tabular}\\
$\Rightarrow$ \quad (quick, the); (quick, brown); (quick, fox)
\begin{tabular}{|c|c|c|c|ccccc|}
\hline
\cellcolor{blue!65}The & \cellcolor{blue!65}quick & \cellcolor{blue!15}brown & \cellcolor{blue!65}fox & \cellcolor{blue!65}jumps & over & the & lazy & dog \\
\hline
\end{tabular}\\
$\Rightarrow$ \quad (brown, the); (brown, quick); (brown, fox); (brown, blue)
\begin{tabular}{|c|c|c|c|ccccc|}
\hline
The & \cellcolor{blue!65}quick & \cellcolor{blue!65}brown & \cellcolor{blue!15}fox & \cellcolor{blue!65}jumps & \cellcolor{blue!65}over & the & lazy & dog \\
\hline
\end{tabular}\\
$\Rightarrow$ \quad (fox, quick); (fox, brown); (fox, jumps); (fox, over)
\end{block}
	
\vfill

\end{vbframe}

% ------------------------------------------------------------------------------

\begin{vbframe}{Self-supervised objectives}

\vfill

\textbf{Self-supervised objectives:}
	
	\begin{itemize}
		\item Skip-gram objective (cf. word2vec \href{https://arxiv.org/pdf/1301.3781.pdf}{\beamergotobutton{Mikolov et al. (2013a)}})
		\item Language modeling objective (cf. \href{http://www.jmlr.org/papers/volume3/bengio03a/bengio03a.pdf}{\beamergotobutton{Bengio et al. (2003)}})
		\item \textit{Masked language modeling (MLM)} objective (cf. chapter 6)\\
					$\rightarrow$ Replace words by a \texttt{[MASK]} token and train the model to predict
		\item \textit{Permutation language modeling (PLM)} objective (cf. chapter 6)\\
					$\rightarrow$ Autoregressive objective of XLNet
		\item \textit{Replaced token detection} objective (cf. chapter 6)\\
					$\rightarrow$ Requires two models: One performing MLM \& the second model to discriminate between actual and the predicted tokens
	\end{itemize}
\vfill

\end{vbframe}

% ------------------------------------------------------------------------------

\begin{vbframe}{Pre-training vs. fine-tuning}

\vfill

\textbf{Pre-training:}

	\begin{itemize}
		\item Using unlabeled corpora in conjunction with self-supervised objectives 
					is commonly referred to as \textit{Pre-Training} the model
		\item Generation of samples for pre-training basically effortless, exploiting 
					the inherent structure of the text
		\item Construction of different self-supervised objectives, which are assumed
			\begin{itemize}
				\item to cover different phenomena better than the others
				\item to work more efficiently for learning
			\end{itemize}
	\end{itemize}
\vfill

\end{vbframe}

% ------------------------------------------------------------------------------

\begin{vbframe}{Pre-training vs. fine-tuning}

\vfill

\textbf{Fine-tuning:}

	\begin{itemize}
		\item The second phase of transfer learning, i.e. adapting the pre-trained model
					to a labeled data set for a specific downstream task is referred to as \textit{Fine-Tuning}
		\item Far less labeled data required compared to a scenario w/o pre-training
		\item Also possible (cf. chapter 7): No fine-tuning, but .. 
			\begin{itemize}
				\item \textit{Zero-Shot Transfer} w/o ANY labeled data
				\item \textit{Few-Shot Transfer} w/ FEW labaled data points
			\end{itemize}
		\item In both of the latter cases, good pre-training becomes even more important
	\end{itemize}
	
\vfill

\end{vbframe}

% ------------------------------------------------------------------------------

\endlecture
\end{document}
